\chapter{\threer}\label{three}
    \section{VirusTotal}
        \subsection{Detection}
        The antimalware analysis with VirusTotal shows that:
        \\\three
        \begin{itemize}
            \item 42 out of 78 security vendors flagged the file as "\textbf{malicious}";
            \item 29 out of 78 security vendors flagged the file as a "\textbf{secure file}";
            \item 1 out of 78 security vendors get a "\textbf{timeout error}";
            \item 7 out of 78 security vendors get a "\textbf{Unable to process file type}" error
        \end{itemize}
        
        Most flags correctly identify the sample as a generic android trojan, but some of them (e.g. CTX,
         Bitdefender, Arcabit) are even able to recognize its malware family (downloader). 
        A few vendors classify the sample as a generic android malware, without specifying the family nor the type (trojan).\\
         It’s interesting to notice that some famous vendor like Microsoft Defender, McAfee and Malwarebytes
         were unable to identify the file as malicious.
        \subsection{Details}
            \todo[inline]{Add more information about Details tab}
        \subsection{Relations}
            \todo[inline]{Add more information about Relations tab}
        \subsection{Behavior}
            \todo[inline]{Add more information about Behavior tab}
    \section{MobSF}
            \todo[inline]{Add MobSF malware analysis}
    \section{Gennymotion}
            \todo[inline]{Add Gennymotion malware analysis}